\documentclass[../Main.tex]{subfiles}
\begin{document}
  \chapter{2022-2023学年微积分（一）（上）期末考试}

  \section{单项选择题(每小题 3 分，共 18 分)}
  \begin{enumerate}
    \item 若$f'(x_{0})=\frac{1}{2}$，则当$\Delta x\to0$时，函数$y=f(x)$在$x=x_{0}$处的微分$\mathrm{d}
      y$是与$\Delta x$（\qquad）的无穷小.

      \begin{tasks}[label=\textbf{\Alph*.}, label-width=1.5em, column-sep=1em](2)
        \task 高阶
        \task 低阶
        \task 同阶
        \task 等价
      \end{tasks}

    \item 设$f(x)$在$x=0$处连续，则$f(x)$在$x=0$处可导的充分条件是（\qquad）.

      \begin{tasks}[label=\textbf{\Alph*.}, label-width=1.5em, column-sep=1em](1)
        \task $\lim_{x\to0}\frac{f(x)-f(-x)}{2x}$存在
        \task $\lim_{x\to0}\frac{f(x^{2})-f(0)}{x^{2}}$存在
        \task $\lim_{x\to0}\frac{f(x)-f(0)}{\sqrt[3]{x}}$存在
        \task $\lim_{x\to\infty}x f\left(\frac{1}{x}\right)$存在
      \end{tasks}

    \item 下列关于数列的描述中，正确的是（\qquad）.

      \begin{tasks}[label=\textbf{\Alph*.}, label-width=1.5em, column-sep=1em](1)
        \task 若$\lim_{n\to\infty}x_{n} y_{n}=0$，则必有$\lim_{n\to\infty}x_{n}=0$或$\lim_{n\to\infty}y_{n}=0$
        \task 若$\{x_{n} y_{n}\}$有界，则必有$\{x_{n}\}$有界或$\{y_{n}\}$有界
        \task 若$\lim_{n\to\infty}(x_{n}-y_{n})=0$，则$\lim_{n\to\infty}x_{n}=\lim_{n\to\infty}y_{n}$
        \task 若有区间$\mathrm{I}$内的数列$x_{n}$，使$|f(x_{n})|$无界，则$f(x)$在$\mathrm{I}$内无界
      \end{tasks}

    \item 设$f(x)$在区间$(a,+\infty)$内可导，且$f'(x)>0$，若$f(a)=0$，则在区间$(a
      ,+\infty)$内有（\qquad）.

      \begin{tasks}[label=\textbf{\Alph*.}, label-width=1.5em, column-sep=1em](2)
        \task $f(x)\geq0$
        \task $f(x)>0$
        \task 不能确定$f(x)$的符号
        \task $f(x)$单调趋向于$+\infty$
      \end{tasks}

    \item 已知$f(x)$在$x=0$处连续，且$\lim_{x\to0}\frac{f(x)}{x^{2}}=1$，则下列结论成立的是（\qquad）.

      \begin{tasks}[label=\textbf{\Alph*.}, label-width=1.5em, column-sep=1em](2)
        \task $f'(0)$存在，且$f'(0)\neq0$
        \task $f'(0)$不存在
        \task $f(x)$在$x=0$处取得极小值
        \task $f(x)$在$x=0$处取得极大值
      \end{tasks}

    \item $\lim_{n\to\infty}\ln\sqrt[n]{\left(1+\frac{1^{2}}{n^{2}}\right)\left(1+\frac{2^{2}}{n^{2}}\right)\cdots\left(1+\frac{n^{2}}{n^{2}}\right)}
      =$（\qquad）.

      \begin{tasks}[label=\textbf{\Alph*.}, label-width=1.5em, column-sep=1em](1)
        \task $\int_{1}^{2}\ln(1+x^{2})\,\mathrm{d}x$
        \task $2\int_{1}^{2}\ln x\,\mathrm{d}x$
        \task $\int_{0}^{1}\ln(1+x^{2})\,\mathrm{d}x$
        \task $2\int_{0}^{1}\ln x\,\mathrm{d}x$
      \end{tasks}
  \end{enumerate}

  \section{填空题(每小题 4 分，共 16 分)}
  \begin{enumerate}
    \setcounter{enumi}{6}
    \item 当$x\to0$时，无穷小量$\left(1+x\right)^{x^2}-1$的阶数是$=\underline
      {\hspace{3cm}}$.
      \vspace{0.5em}

    \item 曲线$\cot\left(x+y+\frac{\pi}{4}\right)=\mathrm{e}^{y}$在点$(0,0)$处的切线方程为$\underline
      {\hspace{3cm}}$.
      \vspace{0.5em}

    \item 若$\int_{0}^{x}f(t)\,\mathrm{d}t=x\mathrm{e}^{-x}$，则$\int_{1}^{+\infty}
      \frac{f(\ln x)}{x}\,\mathrm{d}x=\underline{\hspace{3cm}}$.
      \vspace{0.5em}

    \item 曲线$y=\int_{0}^{x}\tan t\,\mathrm{d}t\left(0\leq x\leq \frac{\pi}{4}
      \right)$的弧长$s=\underline{\hspace{3cm}}$.
      \vspace{0.5em}
  \end{enumerate}

  \section{计算题(每小题 7 分，共 42 分)}
  \begin{enumerate}
    \setcounter{enumi}{10}
    \item 求曲线$C:y=x\arctan x(x>0)$的渐近线.
      \vspace{12em}

    \item 设$f(x)$在$x=a$的某邻域内可导，且$f(a)=f'(a)=1$，求$l=\lim_{x\to a}
      \left(\frac{1}{x-a}-\frac{1}{\displaystyle\int_{a}^{x}f(t)\,\mathrm{d}t}\right
      )$.
      \vspace{11em}

    \item 求$I=\int\frac{1}{x(x^{n}+1)}\,\mathrm{d}x$，其中$n$为正整数.
      \vspace{11em}

    \item 求微分方程$\frac{\mathrm{d}y}{\mathrm{d}x}=\frac{y}{2x+y^{4}}$的通解.
      \vspace{11em}

    \item 设$f(x)=\int_{x}^{1}\sin t^{2}\,\mathrm{d}t$，计算$I=\int_{0}^{1}f(x
      )\,\mathrm{d}x$.
      \vspace{12em}

    \item 求$F(x)=\int_{-\frac{\pi}{2}}^{x}\mathrm{e}^{-t}\sin t\,\mathrm{d}t$在区间$\left
      (-\frac{\pi}{2},\frac{\pi}{2}\right)$内的极值.
      \vspace{12em}
  \end{enumerate}

  \section{综合题(每小题 7 分，共 14 分)}
  \begin{enumerate}
    \setcounter{enumi}{16}
    \item 已知函数$f(x)=\int_{1}^{x}\sqrt{1+t^{2}}\,\mathrm{d}t+\int_{x^2}^{1}
      \sqrt{1+t}\,\mathrm{d}t$，讨论方程$f(x)=0$的实根个数.
      \vspace{12em}

    \item 设$f(x)$满足$\int_{0}^{x}f(t-x)\,\mathrm{d}t=-\frac{x^{3}}{3}-x^{2}$，求曲线$y
      =f(x)$与$x$轴所围成的图形绕$y$轴旋转一周所得的旋转体的体积.
      \vspace{13em}
  \end{enumerate}

  \section{证明题(每小题 5 分，共 10 分)}
  \begin{enumerate}
    \setcounter{enumi}{18}
    \item 设函数$f(x),g(x)$在$[0,1]$上连续，且$f(x)>0$，$g(x)\geq0$.证明：
      \[
        \lim_{n\to\infty}\int_{0}^{1}g(x)\sqrt[n]{f(x)}\,\mathrm{d}x=\int_{0}^{1}g(
        x)\,\mathrm{d}x.
      \]
      \vspace{14em}

    \item 设$f(x)$在$[0,2]$上二阶可导，且$f'(0)=f'(2)=0$. 试证：存在$\xi\in(
      0,2)$，使
      \[
        |f''(\xi)|\geq|f(2)-f(0)|.
      \]
      \vspace{14em}
  \end{enumerate}

  \chapter{2022-2023学年微积分（一）（上）期末考试参考答案}
  \section{单项选择题 (每小题 3 分，共 18 分)}
  \begin{enumerate}
    \item \textbf{Solution}. C.

      由导数的定义，$f'(x_{0})=\lim_{\Delta x\to0}\frac{f(x_{0}+\Delta x)-f(x_{0})}{\Delta
      x}=\lim_{\Delta x\to0}\frac{\mathrm{d}y}{\Delta x}=\frac{1}{2}$.

      因此$\mathrm{d}y$与$\Delta x$是同阶无穷小，但不是等价无穷小.

    \item \textbf{Solution}. D.

      对于A选项，考虑$f(x)=|x|$，则
      \[
        \lim_{x\to0^-}\frac{f(x)-f(-x)}{2x}=\lim_{x\to0^-}\frac{(-x)-(-x)}{2x}= 0
        \text{，}\lim_{x\to0^+}\frac{f(x)-f(-x)}{2x}=\lim_{x\to0^+}\frac{x-x}{2x}
        =0,
      \]

      故$\lim_{x\to0}\frac{f(x)-f(-x)}{2x}=0$存在，但显然$f(x)$在$x=0$处不可导.

      对于B选项，令$t=x^{2}$，则$\lim_{x\to0}\frac{f(x^{2})-f(0)}{x^{2}}=\lim_{t\to0^+}
      \frac{f(t)-f(0)}{t}$存在，仅能说明$f'_{+}(0)$存在.

      对于C选项，考虑$f(x)=\sqrt[3]{x}$，则$\lim_{x\to0}\frac{f(x)-f(0)}{\sqrt[3]{x}}
      =1$存在，但显然$f(x)$在$x=0$处不可导.

      对于D选项，令$t=\frac{1}{x}$，则$\lim_{x\to\infty}x f\left(\frac{1}{x}\right
      )=\lim_{t\to0}\frac{f(t)}{t}$存在，

      结合$f(x)$在$x=0$处的连续性可得$f(0)=\lim_{t\to0}f(t)=\lim_{t\to0}\frac{f(t)}{t}
      \cdot t=0$，

      所以$\lim_{t\to0}\frac{f(t)-f(0)}{t}=\lim_{t\to0}\frac{f(t)}{t}$存在，即$f(
      x)$在$x=0$处可导.

    \item \textbf{Solution}. D.

      对于A，B，C选项，考虑$x_{n}=
      \begin{cases}
        0, & n=2k-1, \\
        n, & n=2k
      \end{cases}$，$y_{n}=
      \begin{cases}
        n, & n=2k-1, \\
        0, & n=2k
      \end{cases}$，其中$k$是正整数.

      则$x_{n} y_{n}\equiv0$，但$x_{n}$，$y_{n}$均无界，且$\lim_{n\to\infty}x_{n}$，$\lim
      _{n\to\infty}y_{n}$均不存在.

    \item \textbf{Solution}. C.

      假设$a>0$，考虑函数$f(x)=
      \begin{cases}
        -\frac{1}{x}+\frac{1}{2a}, & x>a, \\
        0,                         & x=a.
      \end{cases}$

      显然$f(x)$在区间$(a,+\infty)$内可导，当$a<x<2a$时，$f(x)<0$；当$x>2a$时，$f
      (x)>0$，

      且$\lim_{x\to+\infty}f(x)=\frac{1}{2a}<+\infty$.

    \item \textbf{Solution}. C.

      结合$f(x)$在$x=0$处的连续性可得$f(0)=\lim_{x\to0}f(x)=\lim_{x\to0}\frac{f(x)}{x^{2}}
      \cdot x^{2}=0$，

      因此$\lim_{x\to0}\frac{f(x)-f(0)}{x}=\lim_{x\to0}\frac{f(x)}{x^{2}}\cdot x=
      0$，即$f'(0)=0$.

      又由极限的保号性，当$x\to0$时，$\frac{f(x)}{x^{2}}>0$，所以$f(x)>0$. 即$x=0$是$f
      (x)$的极小值点.

    \item \textbf{Solution}. C.

      由定积分的定义可得
      \[
        \begin{aligned}
          \lim_{n\to\infty}\ln\sqrt[n]{\left(1+\frac{1^2}{n^2}\right)\left(1+\frac{2^2}{n^2}\right)\cdots\left(1+\frac{n^2}{n^2}\right)}= & \lim_{n\to\infty}\frac{1}{n}\sum_{k=1}^{n}\ln\left(1+\frac{k^2}{n^2}\right) \\
          =                                                                                                                               & \int_{0}^{1}\ln(1+x^{2})\,\mathrm{d}x.
        \end{aligned}
      \]
  \end{enumerate}

  \section{填空题(每小题 4 分，共 16 分)}
  \begin{enumerate}
    \setcounter{enumi}{6}
    \item \textbf{Solution}. $3$.

      当$x\to0$时，$\left(1+x\right)^{x^2}-1 = \mathrm{e}^{x^2\ln(1+x)}-1 = \mathrm{e}
      ^{x^2\left(x+o(x)\right)}-1= \mathrm{e}^{x^3+o(x^3)}-1 = x^{3}+o(x^{3})$.

      所以无穷小量$\left(1+x\right)^{x^2}-1$的阶数是$3$.

    \item \textbf{Solution}. $2x+3y=0$.

      方程$\cot\left(x+y+\frac{\pi}{4}\right)=\mathrm{e}^{y}$两边对$x$求导得
      \[
        -\csc^{2}\left(x+y+\frac{\pi}{4}\right)\cdot\left(1+y'\right)=\mathrm{e}^{y}
        y'.
      \]

      将$x=0$，$y=0$代入上式解得$y'(0)=-\frac{2}{3}$，所以切线方程为$2x+3y=0$.

    \item \textbf{Solution}. $0$.

      由题可知，$f(x)=\left(x\mathrm{e}^{-x}\right)'=\mathrm{e}^{-x}-x\mathrm{e}^{-x}
      =(1-x)\mathrm{e}^{-x}$. 故$f(\ln x)=\frac{1-\ln x}{x}$.

      所以
      \[
        \begin{aligned}
          \int_{1}^{+\infty}\frac{f(\ln x)}{x}\,\mathrm{d}x = & \int_{1}^{+\infty}\frac{1-\ln x}{x^2}\,\mathrm{d}x                                                                                                                  \\
          =                                                 & \int_{1}^{+\infty}\left(\ln x-1\right)\,\mathrm{d}\left(\frac{1}{x}\right)= \left.\frac{\ln x-1}{x}\right|_{1}^{+\infty}-\int_{1}^{+\infty}\frac{1}{x^2}\,\mathrm{d}x \\
          =                                                 & 1+\left.\frac{1}{x}\right|_{1}^{+\infty}= 0.
        \end{aligned}
      \]

    \item \textbf{Solution}. $\ln\left(1+\sqrt{2}\right)$.

      $y'=\tan x$，所以$\mathrm{d}s=\sqrt{1+y'^{2}}\,\mathrm{d}x=\sqrt{1+\tan^{2} x}
      \mathrm{d}x=\sec x\,\mathrm{d}x$，

      故$s=\int_{0}^{\frac{\pi}{4}}\sec x\,\mathrm{d}x=\ln\left|\sec x+\tan x\right
      |\bigg|_{0}^{\frac{\pi}{4}}=\ln\left(1+\sqrt{2}\right)$.
  \end{enumerate}
  \section{计算题(每小题 7 分，共 42 分)}
  \begin{enumerate}
    \setcounter{enumi}{10}
    \item \textbf{Solution}. 曲线$C$没有竖直渐近线和水平渐近线.

      $\lim_{x\to+\infty}\frac{y}{x}= \lim_{x\to+\infty}\arctan x = \frac{\pi}{2}$，

      $\lim_{x\to+\infty}\left(y-\frac{\pi}{2}x\right)=\lim_{x\to+\infty}x\left(\arctan
      x-\frac{\pi}{2}\right)=\lim_{x\to+\infty}\frac{\arctan x-\frac{\pi}{2}}{\frac{1}{x}}
      =-\lim_{x\to+\infty}\frac{x^{2}}{1+x^{2}}=-1$.

      所以$y=\frac{\pi}{2}x-1$是曲线$C$的斜渐近线.

    \item \textbf{Solution}. 由积分中值定理，存在$\xi$介于$a$和$x$之间，使得$\int
      _{a}^{x}f(t)\,\mathrm{d}t=f(\xi)(x-a)$，

      当$x\to a$时，$\xi\to a$，所以
      \[
        \begin{aligned}
          l = & \lim_{x\to a}\frac{\displaystyle\int_{a}^{x}f(t)\,\mathrm{d}t-(x-a)}{(x-a)\displaystyle\int_{a}^{x}f(t)\,\mathrm{d}t}      \\
          =   & \lim_{x\to a}\frac{\displaystyle\int_{a}^{x}f(t)\,\mathrm{d}t-(x-a)}{(x-a)^2 f(\xi)}= \lim_{x\to a}\frac{f(x)-1}{2(x-a)} \\
          =   & \frac{1}{2}\lim_{x\to a}\frac{f(x)-f(a)}{x-a}=\frac{1}{2}f'(a)=\frac{1}{2}.
        \end{aligned}
      \]

    \item \textbf{Solution}.
      \[
        \begin{aligned}
          I = \int\frac{1}{x(x^n+1)}\,\mathrm{d}x = & \int\frac{x^{-n-1}}{x^{-n}+1}\,\mathrm{d}x                          \\
          =                                       & -\frac{1}{n}\int\frac{1}{x^{-n}+1}\,\mathrm{d}\left(x^{-n}+1\right) \\
          =                                       & -\frac{1}{n}\ln\left|x^{-n}+1\right|+C.
        \end{aligned}
      \]

    \item \textbf{Solution}. 方程变形为$\frac{\mathrm{d}x}{\mathrm{d}y}-\frac{2x}{y}
      =y^{3}$.

      由一阶非齐次线性微分方程的通解公式得
      \[
        \begin{aligned}
          x = & \mathrm{e}^{-\int\left(-\frac{2}{y}\right)\,\mathrm{d}y}\left(C+\int y^{3}\mathrm{e}^{\int\left(-\frac{2}{y}\right)\,\mathrm{d}y}\,\mathrm{d}y\right) \\
          =   & y^{2}\left(C+\int y\,\mathrm{d}y\right) = y^{2}\left(C+\frac{1}{2}y^{2}\right) = Cy^{2}+\frac{1}{2}y^{4}.
        \end{aligned}
      \]

      其中C为任意常数.

    \item \textbf{Solution}. 由题可知$f(1)=0$，$f'(x)=-\sin x^{2}$，所以
      \[
        \begin{aligned}
          I = \int_{0}^{1}f(x)\,\mathrm{d}x = & xf(x)\bigg|_{0}^{1}-\int_{0}^{1}x f'(x)\,\mathrm{d}x                                    \\
          =                                 & \int_{0}^{1}x\sin x^{2}\,\mathrm{d}x = \frac{1}{2}\int_{0}^{1}\sin x^{2}\,\mathrm{d}x^{2} \\
          =                                 & -\frac{1}{2}\cos x^{2}\bigg|_{0}^{1}                                                  \\
          =                                 & \frac{1-\cos 1}{2}.
        \end{aligned}
      \]

    \item \textbf{Solution}. $F'(x)=\mathrm{e}^{-x}\sin x$，在$\left(-\frac{\pi}{2}
      ,\frac{\pi}{2}\right)$内，令$F'(x)=0$得$x=0$.

      计算
      \[
        \begin{aligned}
          \int \mathrm{e}^{-t}\sin t\,\mathrm{d}t = & -\int \mathrm{e}^{-t}\,\mathrm{d}\cos t = -\mathrm{e}^{-t}\cos t - \int \mathrm{e}^{-t}\cos t\,\mathrm{d}t \\
          =                                       & -\mathrm{e}^{-t}\cos t - \int \mathrm{e}^{-t}\,\mathrm{d}\sin t                                          \\
          =                                       & -\mathrm{e}^{-t}\cos t - \mathrm{e}^{-t}\sin t - \int \mathrm{e}^{-t}\sin t\,\mathrm{d}t
        \end{aligned}
      \]

      所以$\int \mathrm{e}^{-t}\sin t\,\mathrm{d}t=-\frac{1}{2}\mathrm{e}^{-t}(\sin
      t+\cos t)+C$，

      因此极值$F(0)=\left.-\frac{1}{2}\mathrm{e}^{-t}(\sin t+\cos t)\right|_{-\frac{\pi}{2}}
      ^{0}=-\frac{1}{2}\left(\mathrm{e}^{\frac{\pi}{2}}+1\right)$.
  \end{enumerate}

  \section{综合题(每小题 7 分，共 14 分)}
  \begin{enumerate}
    \setcounter{enumi}{16}
    \item \textbf{Solution}. $f'(x)=\sqrt{1+x^{2}}-2x\sqrt{1+x^{2}}=(1-2x)\sqrt{1+x^{2}}$，

      所以$f(x)$在$\left(-\infty,\frac{1}{2}\right)$上单调增加，在$\left(\frac{1}{2}
      ,+\infty\right)$上单调减少.

      因为$f(0)=\int_{1}^{0}\sqrt{1+t^{2}}\,\mathrm{d}t+\int_{0}^{1}\sqrt{1+t}\,\mathrm{d}
      t=\int_{0}^{1}\left(\sqrt{1+t}-\sqrt{1+t^{2}}\right)\,\mathrm{d}t>0$，

      $f(1)=0$，$f(-1)=\int_{1}^{-1}\sqrt{1+t^{2}}\,\mathrm{d}t<0$，故方程$f(x)=0$有两个实根.

    \item \textbf{Solution}. 令$t-x=u$，则$\int_{-x}^{0}f(u)\,\mathrm{d}u=-\frac{x^{3}}{3}
      -x^{2}$.

      方程两边求导得$f(-x)=-x^{2}-2x$，所以$f(x)=-x^{2}+2x$.

      所以旋转体的体积$V=$
      \[
        \begin{aligned}
          \int_{0}^{2}2\pi xf(x)\,\mathrm{d}x = & 2\pi\int_{0}^{2}x(-x^{2}+2x)\,\mathrm{d}x = 2\pi\int_{0}^{2}(-x^{3}+2x^{2})\,\mathrm{d}x \\
          =                                   & 2\pi\left[-\frac{x^4}{4}+\frac{2x^3}{3}\right]_{0}^{2}= \frac{8}{3}\pi.
        \end{aligned}
      \]
  \end{enumerate}

  \section{证明题(每小题 5 分，共 10 分)}
  \begin{enumerate}
    \setcounter{enumi}{18}
    \item \textbf{Proof}. 易知函数$\sqrt[n]{f(x)}$在闭区间$[0,1]$上连续，因而具有最小值$m$和最大值$M$.所以
      \[
        m\leq \frac{\displaystyle\int_{0}^{1}g(x)\sqrt[n]{f(x)}\,\mathrm{d}x}{\displaystyle\int_{0}^{1}g(x)\,\mathrm{d}x}
        \leq M.
      \]

      由介值定理，存在$\xi\in[0,1]$使得$\sqrt[n]{f(\xi)}=\frac{\displaystyle\int_{0}^{1}g(x)\sqrt[n]{f(x)}\,\mathrm{d}x}{\displaystyle\int_{0}^{1}g(x)\,\mathrm{d}x}$.

      上式两边取极限，由于$\lim_{n\to\infty}\sqrt[n]{f(\xi)}=1$，所以$\lim_{n\to\infty}
      \frac{\displaystyle\int_{0}^{1}g(x)\sqrt[n]{f(x)}\,\mathrm{d}x}{\displaystyle\int_{0}^{1}g(x)\,\mathrm{d}x}
      =1$，

      即$\lim_{n\to\infty}\int_{0}^{1}g(x)\sqrt[n]{f(x)}\,\mathrm{d}x=\int_{0}^{1}g
      (x)\,\mathrm{d}x$.

    \item \textbf{Proof}. 由Taylor公式，将$f(x)$分别在$x=0$和$x=2$处展开得
      \[
        \begin{gathered}
          f(x)=f(0)+f'(0)x+\frac{1}{2}f''(\eta)x^2=f(0)+\frac{1}{2}f''(\eta)x^2,\\
          f(x)=f(2)+f'(2)(x-2)+\frac{1}{2}f''(\zeta)(x-2)^2=f(2)+\frac{1}{2}f''(\zeta)(x-2)^2,
        \end{gathered}
      \]

      两式相减得
      \[
        f(2)-f(0)=\frac{1}{2}\left[f''(\zeta)(x-2)^{2}-f''(\eta)x^{2}\right].
      \]

      所以
      \[
        |f(2)-f(0)|\leq\frac{1}{2}\left[\left|f''(\zeta)\right|(x-2)^{2}+\left|f''(\eta)\right|x^{2}\right].
      \]

      取$|f''(\xi)|=\max\{|f''(\eta)|,|f''(\zeta)|\}$，

      则$|f(2)-f(0)|\leq\frac{1}{2}|f''(\xi)|\left[(x-2)^{2}+x^{2}\right]\leq \frac{1}{2}
      |f''(\xi)|\cdot 2=|f''(\xi)|$，即$|f''(\xi)|\geq|f(2)-f(0)|$.
  \end{enumerate}
\end{document}




